% !TEX program = xelatex
\documentclass[12pt,a4paper]{article}

% ── 中文支持 ──
\usepackage{ctex}
\setCJKmainfont{SimSun}[BoldFont=SimHei,ItalicFont=KaiTi]
\setCJKsansfont{Microsoft YaHei}
\setCJKmonofont{FangSong}

% ── 页面与排版 ──
\usepackage[top=2.5cm,bottom=2.5cm,left=2.8cm,right=2.8cm]{geometry}
\usepackage{setspace}
\onehalfspacing
\usepackage{titlesec}
\titleformat{\section}{\Large\bfseries\sffamily}{\thesection}{1em}{}
\titleformat{\subsection}{\large\bfseries\sffamily}{\thesubsection}{1em}{}
\titleformat{\subsubsection}{\normalsize\bfseries\sffamily}{\thesubsubsection}{1em}{}

% ── 图表与颜色 ──
\usepackage{graphicx}
\usepackage{float}
\usepackage{booktabs}
\usepackage{array}
\usepackage{longtable}
\usepackage{multirow}
\usepackage{xcolor}
\definecolor{primary}{HTML}{1E3A5F}
\definecolor{accent}{HTML}{2B6CB0}
\definecolor{lightbg}{HTML}{F0F4F8}
\definecolor{codebg}{HTML}{F5F5F5}

% ── 代码 ──
\usepackage{listings}
\lstset{
  basicstyle=\small\ttfamily,
  backgroundcolor=\color{codebg},
  frame=single,
  framerule=0.5pt,
  rulecolor=\color{gray!40},
  breaklines=true,
  columns=fullflexible,
  keepspaces=true,
  tabsize=4,
  showstringspaces=false,
  keywordstyle=\color{accent}\bfseries,
  commentstyle=\color{gray},
  stringstyle=\color{teal},
  numbers=left,
  numberstyle=\tiny\color{gray},
  numbersep=8pt,
  xleftmargin=1.5em,
}

% ── 超链接与引用 ──
\usepackage{hyperref}
\hypersetup{
  colorlinks=true,
  linkcolor=primary,
  urlcolor=accent,
  citecolor=primary,
  pdfauthor={DCF Valuation Agent Team},
  pdftitle={DCF估值分析智能代理——项目技术报告},
}

% ── 数学公式 ──
\usepackage{amsmath,amssymb}

% ── 页眉页脚 ──
\usepackage{fancyhdr}
\setlength{\headheight}{14.5pt}
\pagestyle{fancy}
\fancyhf{}
\fancyhead[L]{\small\sffamily DCF估值分析智能代理}
\fancyhead[R]{\small\sffamily 项目技术报告}
\fancyfoot[C]{\thepage}
\renewcommand{\headrulewidth}{0.4pt}

% ── 目录样式 ──
\usepackage{tocloft}
\renewcommand{\cftsecleader}{\cftdotfill{\cftdotsep}}

% ── 引用 ──
\usepackage{enumitem}

% ── 盒子 ──
\usepackage{tcolorbox}
\tcbuselibrary{skins,breakable}
\newtcolorbox{keybox}[1][]{
  colback=lightbg,
  colframe=accent,
  fonttitle=\bfseries\sffamily,
  title=#1,
  boxrule=0.8pt,
  arc=3pt,
  breakable,
}

% ══════════════════════════════════════════
\begin{document}

% ── 封面 ──
\begin{titlepage}
\begin{center}
\vspace*{3cm}

{\color{primary}\rule{\textwidth}{2pt}}

\vspace{1.5cm}

{\fontsize{28}{34}\selectfont\sffamily\bfseries\color{primary}
DCF 估值分析智能代理}

\vspace{0.8cm}

{\fontsize{16}{20}\selectfont\sffamily\color{accent}
基于大语言模型的贴现现金流自动估值系统}

\vspace{0.5cm}

{\color{primary}\rule{0.6\textwidth}{1pt}}

\vspace{2cm}

{\large\sffamily 项目技术报告}

\vspace{3cm}

\begin{tabular}{rl}
\textbf{LLM 模型：} & MiniMax-M2.7-highspeed \\[6pt]
\textbf{后端框架：} & Python FastAPI \\[6pt]
\textbf{前端框架：} & React 18 + TypeScript + Vite \\[6pt]
\textbf{日期：} & 2026 年 3 月 \\
\end{tabular}

\vfill

{\color{primary}\rule{\textwidth}{2pt}}
\end{center}
\end{titlepage}

% ── 目录 ──
\tableofcontents
\newpage

% ══════════════════════════════════════════
\section{项目概述}

\subsection{背景与动机}

随着金融市场数据量的持续激增和财务分析复杂度的不断提升，传统的 DCF（Discounted Cash Flow，贴现现金流）估值方法面临显著的效率瓶颈。传统方法依赖大量人工操作：分析师需要手动从上市公司年报中提取关键财务数据、逐项整理 Excel 模型、反复调整假设参数，整个过程耗时长且容易出错。

本项目旨在设计并实现一个\textbf{基于大语言模型与自动化工作流的 DCF 估值智能代理}，实现从企业财报上传到最终估值结果的全自动化流程，为投资分析师和金融研究人员提供高效、准确的估值支持工具。

\subsection{核心目标}

\begin{enumerate}[leftmargin=2em]
  \item \textbf{自动化数据提取}：上传 PDF 格式的上市公司年报，由 AI 自动提取 20+ 项关键财务指标。
  \item \textbf{完整 DCF 建模}：基于 FCFF 方法构建自由现金流折现模型，计算企业价值与每股内在价值。
  \item \textbf{多维敏感性分析}：生成 WACC 与永续增长率的交叉敏感性矩阵，量化参数假设对估值的影响。
  \item \textbf{AI 估值报告}：利用大语言模型自动生成专业级估值分析叙述。
  \item \textbf{美观的交互界面}：提供完整的前后端系统，支持中英文双语切换与深色/浅色主题。
\end{enumerate}

\subsection{技术选型概览}

\begin{table}[H]
\centering
\caption{系统技术栈}
\label{tab:techstack}
\begin{tabular}{@{}lll@{}}
\toprule
\textbf{层级} & \textbf{技术} & \textbf{说明} \\
\midrule
后端框架  & Python FastAPI + Uvicorn & 高性能异步 Web 框架 \\
数据模型  & Pydantic v2            & 类型安全的请求/响应校验 \\
PDF 解析  & pdfplumber             & 精确提取 PDF 文本与表格 \\
LLM 服务  & MiniMax-M2.7-highspeed & 推理型大语言模型 \\
LLM SDK   & OpenAI Python SDK      & 兼容 MiniMax API 端点 \\
前端框架  & React 18 + TypeScript  & 组件化前端开发 \\
构建工具  & Vite 5                 & 极速热更新开发体验 \\
UI 组件库 & Ant Design 5           & 企业级 React 组件 \\
样式方案  & TailwindCSS            & 原子化 CSS 实用类 \\
图表库    & ECharts                & 专业级交互式图表 \\
状态管理  & Zustand                & 轻量级全局状态 \\
国际化    & react-i18next          & 中英文双语切换 \\
\bottomrule
\end{tabular}
\end{table}

% ══════════════════════════════════════════
\section{系统架构设计}

\subsection{整体架构}

系统采用经典的\textbf{前后端分离}架构，前端通过 Vite 开发服务器代理将 API 请求转发至后端 FastAPI 服务。整体数据流如下：

\begin{keybox}[系统数据流]
\begin{enumerate}[leftmargin=2em]
  \item 用户上传 PDF 财务报告或手动输入财务数据；
  \item 后端使用 pdfplumber 智能提取 PDF 关键页面文本；
  \item 提取的文本发送至 MiniMax LLM，返回结构化 JSON 财务数据；
  \item 用户可在界面中审查和编辑提取结果，并调整 DCF 参数；
  \item 后端 DCF 引擎执行估值计算、敏感性分析；
  \item LLM 生成专业估值叙述报告；
  \item 前端以图表和仪表盘形式展示完整估值结果。
\end{enumerate}
\end{keybox}

\subsection{后端架构}

后端采用\textbf{分层架构}，清晰划分 API 层、服务层和模型层：

\begin{table}[H]
\centering
\caption{后端模块划分}
\begin{tabular}{@{}lll@{}}
\toprule
\textbf{模块} & \textbf{文件} & \textbf{职责} \\
\midrule
API 层       & \texttt{api/upload.py}      & PDF 上传与数据提取端点 \\
             & \texttt{api/analysis.py}    & DCF 计算与敏感性分析端点 \\
             & \texttt{api/report.py}      & 报告导出端点（预留） \\
\midrule
服务层       & \texttt{services/pdf\_service.py}       & PDF 智能文本提取 \\
             & \texttt{services/llm\_service.py}       & MiniMax LLM 封装 \\
             & \texttt{services/extractor\_service.py} & LLM 财务数据提取 \\
             & \texttt{services/dcf\_service.py}       & DCF 计算引擎 \\
\midrule
模型层       & \texttt{models/schemas.py}  & Pydantic 数据模型定义 \\
\midrule
配置层       & \texttt{config.py}          & 环境变量与配置加载 \\
             & \texttt{main.py}            & FastAPI 应用入口 \\
\bottomrule
\end{tabular}
\end{table}

\subsection{前端架构}

前端基于 React 单页应用（SPA），使用 React Router 实现页面路由，Zustand 管理全局状态：

\begin{table}[H]
\centering
\caption{前端页面与组件}
\begin{tabular}{@{}lll@{}}
\toprule
\textbf{类型} & \textbf{组件} & \textbf{功能} \\
\midrule
页面 & \texttt{HomePage}     & 首页与引导入口 \\
     & \texttt{AnalysisPage} & 数据输入与参数调整 \\
     & \texttt{ResultPage}   & 估值结果仪表盘 \\
\midrule
组件 & \texttt{FileUpload}        & PDF 拖拽上传 \\
     & \texttt{ManualInput}       & 手动数据输入表单 \\
     & \texttt{ParamPanel}        & DCF 参数滑块面板 \\
     & \texttt{CashFlowChart}     & FCF 预测柱状图 \\
     & \texttt{SensitivityTable}  & 敏感性分析热力图 \\
     & \texttt{WaterfallChart}    & 估值瀑布图 \\
\bottomrule
\end{tabular}
\end{table}

\subsection{API 接口设计}

\begin{table}[H]
\centering
\caption{RESTful API 接口一览}
\label{tab:api}
\begin{tabular}{@{}llp{7cm}@{}}
\toprule
\textbf{方法} & \textbf{路径} & \textbf{说明} \\
\midrule
GET  & \texttt{/api/health}          & 健康检查 \\
POST & \texttt{/api/extract/upload}  & 上传 PDF，提取文本并调用 LLM 返回结构化财务数据 \\
POST & \texttt{/api/extract/text}    & 从纯文本中调用 LLM 提取财务数据 \\
POST & \texttt{/api/calculate}       & 执行 DCF 估值计算，返回完整估值结果 \\
POST & \texttt{/api/sensitivity}     & 生成 WACC-增长率敏感性分析矩阵 \\
POST & \texttt{/api/narrative}       & 调用 LLM 生成专业估值叙述报告 \\
\bottomrule
\end{tabular}
\end{table}

% ══════════════════════════════════════════
\section{核心算法与模型}

\subsection{DCF 估值理论基础}

DCF（Discounted Cash Flow）估值法的核心思想是：\textbf{一项资产的价值等于其未来产生的所有自由现金流的折现值之和}。本系统采用 FCFF（Free Cash Flow to Firm，企业自由现金流）方法进行估值。

\subsection{FCFF 计算}

企业自由现金流的计算公式为：

\begin{equation}
\text{FCFF} = \text{NOPAT} + \text{D\&A} - \text{CapEx} - \Delta\text{NWC}
\label{eq:fcff}
\end{equation}

其中：
\begin{itemize}[leftmargin=2em]
  \item $\text{NOPAT} = \text{EBIT} \times (1 - T_c)$，即税后净营业利润；
  \item $\text{D\&A}$ 为折旧与摊销；
  \item $\text{CapEx}$ 为资本性支出；
  \item $\Delta\text{NWC}$ 为营运资本变动。
\end{itemize}

在预测期内，系统基于用户设定的营收增长率、经营利润率等参数逐年投射 FCFF：

\begin{equation}
\text{Revenue}_t = \text{Revenue}_{t-1} \times (1 + g_{\text{rev}})
\end{equation}

\begin{equation}
\text{FCFF}_t = \text{Revenue}_t \times m_{\text{op}} \times (1 - T_c) + \text{Revenue}_t \times r_{\text{DA}} - \text{Revenue}_t \times r_{\text{CapEx}} - \text{Revenue}_t \times r_{\text{NWC}}
\end{equation}

\subsection{WACC 计算}

加权平均资本成本（WACC）是企业融资成本的加权平均，用作折现率：

\begin{equation}
\text{WACC} = \frac{E}{E+D} \times R_e + \frac{D}{E+D} \times R_d \times (1 - T_c)
\label{eq:wacc}
\end{equation}

其中权益成本 $R_e$ 通过 CAPM（资本资产定价模型）计算：

\begin{equation}
R_e = R_f + \beta \times (R_m - R_f)
\label{eq:capm}
\end{equation}

\begin{itemize}[leftmargin=2em]
  \item $R_f$：无风险利率（默认取中国十年期国债收益率 2.5\%）
  \item $\beta$：股票贝塔系数
  \item $R_m$：市场预期回报率（默认 9\%）
  \item $R_d$：债务成本
  \item $E/(E+D)$：权益占总资本比重
  \item $D/(E+D)$：债务占总资本比重
\end{itemize}

\subsection{终端价值}

预测期结束后，采用 Gordon 永续增长模型计算终端价值（Terminal Value）：

\begin{equation}
\text{TV} = \frac{\text{FCFF}_n \times (1 + g_\infty)}{\text{WACC} - g_\infty}
\label{eq:tv}
\end{equation}

其中 $g_\infty$ 为永续增长率（默认 2.5\%），约束条件为 $\text{WACC} > g_\infty$。

\subsection{企业价值与每股价值}

\begin{equation}
\text{EV} = \sum_{t=1}^{n} \frac{\text{FCFF}_t}{(1 + \text{WACC})^t} + \frac{\text{TV}}{(1 + \text{WACC})^n}
\label{eq:ev}
\end{equation}

\begin{equation}
\text{Equity Value} = \text{EV} - \text{Total Debt} + \text{Cash}
\end{equation}

\begin{equation}
\text{Per Share Value} = \frac{\text{Equity Value}}{\text{Shares Outstanding}}
\end{equation}

\subsection{敏感性分析}

系统构建以 WACC 和永续增长率 $g_\infty$ 为双轴的敏感性矩阵。对于基准值附近 $\pm 2\%$ 范围内的 9 个 WACC 值和 7 个增长率值，分别计算每股价值，生成 $9 \times 7$ 的矩阵。该矩阵以热力图形式可视化，帮助投资者直观理解关键假设对估值的影响程度。

% ══════════════════════════════════════════
\section{LLM 智能提取模块}

\subsection{模型选型}

本系统选用 \textbf{MiniMax-M2.7-highspeed} 作为大语言模型，该模型具有以下特点：

\begin{itemize}[leftmargin=2em]
  \item 支持 OpenAI 兼容 API，可直接使用 OpenAI Python SDK 调用；
  \item 内置推理能力（Chain-of-Thought），适合复杂财务数据解析；
  \item 高速推理，满足实时交互需求；
  \item 原生支持中英文，适合中国 A 股上市公司年报处理。
\end{itemize}

通过 OpenAI SDK 的 \texttt{base\_url} 参数连接 MiniMax API：

\begin{lstlisting}[language=Python,caption=MiniMax LLM 客户端初始化]
from openai import AsyncOpenAI

client = AsyncOpenAI(
    api_key=MINIMAX_API_KEY,
    base_url="https://api.minimaxi.com/v1",
)
\end{lstlisting}

\subsection{PDF 智能提取策略}

面对 200+ 页的中国上市公司年报，直接将全文发送给 LLM 既不经济也超出上下文限制。系统采用\textbf{关键词权重评分}策略：

\begin{enumerate}[leftmargin=2em]
  \item 遍历 PDF 每一页，提取文本内容；
  \item 对每页文本计算关键词匹配得分（关键词包括``营业收入''、``净利润''、``资产负债表''、``现金流''等 17 个中英文财务术语）；
  \item 按得分从高到低排序，优先选取高相关性页面；
  \item 累计字符数不超过 30,000 字（约涵盖 15--20 页高价值内容）；
  \item 按原始页码重新排序后拼接，送入 LLM 进行提取。
\end{enumerate}

\subsection{结构化提取 Prompt 设计}

提取 Prompt 使用中文，精确定义 21 个目标字段及其含义，并明确约定：

\begin{itemize}[leftmargin=2em]
  \item 所有金额统一为``元''（自动处理万元/亿元转换）；
  \item 所有比率使用小数（如 14.23\% $\to$ 0.1423）；
  \item 缺失值基于其他已知数据合理估算；
  \item 仅返回 JSON 对象，不包含 markdown 标记。
\end{itemize}

\subsection{推理标签处理}

MiniMax-M2.7 内置推理模式，会在响应中输出 \texttt{<think>...</think>} 思考过程。系统通过正则表达式自动剥离推理标签，再从剩余文本中定位并解析 JSON 对象：

\begin{lstlisting}[language=Python,caption=推理标签处理]
import re, json

def _extract_json(raw: str) -> dict:
    text = raw.strip()
    text = re.sub(r"<think>.*?</think>", "", text,
                  flags=re.DOTALL).strip()
    match = re.search(r"\{[\s\S]*\}", text)
    if match:
        return json.loads(match.group())
    return json.loads(text)
\end{lstlisting}

% ══════════════════════════════════════════
\section{前端交互设计}

\subsection{页面结构}

前端共包含三个主要页面，遵循\textbf{向导式}交互流程：

\begin{description}[leftmargin=2em]
  \item[首页（HomePage）] 项目介绍、功能亮点展示与``开始分析''入口按钮。
  \item[分析页（AnalysisPage）] 左侧提供``上传报告''与``手动输入''两个标签页，右侧为 DCF 参数实时调整面板（含 WACC、增长率、利润率等滑块），底部为``计算估值''按钮。
  \item[结果页（ResultPage）] 展示估值概要指标、FCF 预测柱状图、敏感性分析热力图、估值瀑布图以及 AI 生成的估值叙述。
\end{description}

\subsection{主要交互功能}

\begin{itemize}[leftmargin=2em]
  \item \textbf{PDF 拖拽上传}：支持点击或拖拽上传 PDF 文件，自动触发数据提取。
  \item \textbf{实时数据编辑}：AI 提取的财务数据以可编辑表单展示，用户可逐项修正。
  \item \textbf{参数滑块调节}：WACC、营收增长率、经营利润率等参数通过滑块和输入框实时调整。
  \item \textbf{一键计算}：点击按钮依次执行 DCF 计算、敏感性分析、AI 叙述生成，完成后自动跳转至结果页。
  \item \textbf{双语切换}：导航栏提供中文/英文一键切换，所有界面文本实时更新。
  \item \textbf{主题切换}：支持浅色/深色模式切换，适应不同工作环境。
\end{itemize}

\subsection{可视化图表}

系统使用 ECharts 库实现三类专业级金融图表：

\begin{enumerate}[leftmargin=2em]
  \item \textbf{FCF 预测柱状图}：展示 5 年预测期内的自由现金流趋势。
  \item \textbf{敏感性热力图}：以 WACC 为纵轴、永续增长率为横轴的 $9 \times 7$ 矩阵，颜色深浅表示每股价值大小。
  \item \textbf{估值瀑布图}：从企业价值出发，经过减除债务、加回现金，最终得到每股价值的桥接图。
\end{enumerate}

% ══════════════════════════════════════════
\section{测试验证}

\subsection{测试样本}

使用\textbf{苏州天华新能源科技股份有限公司（300390）2025 年年度报告}（226 页，约 24.6 万字）作为测试样本，验证系统全流程功能。

\subsection{数据提取结果}

系统成功从 PDF 中自动提取以下关键财务数据（表~\ref{tab:extract}）：

\begin{table}[H]
\centering
\caption{LLM 自动提取结果}
\label{tab:extract}
\begin{tabular}{@{}lrl@{}}
\toprule
\textbf{指标} & \textbf{提取值} & \textbf{单位} \\
\midrule
公司名称       & \multicolumn{2}{l}{苏州天华新能源科技股份有限公司} \\
股票代码       & 300390              & --- \\
报告年度       & 2025                & 年 \\
营业收入       & 7,548,826,104.54    & 元 \\
营收增长率     & 14.23               & \% \\
净利润         & 402,189,233.13      & 元 \\
经营利润率     & 5.33                & \% \\
总债务         & 2,349,172,637.01    & 元 \\
货币资金       & 4,289,026,296.41    & 元 \\
总股本         & 830,750,788         & 股 \\
有效税率       & 1.42                & \% \\
Beta 系数      & 1.0                 & --- \\
\bottomrule
\end{tabular}
\end{table}

\subsection{DCF 估值计算结果}

基于提取数据与默认 DCF 参数（WACC=10\%，营收增长率=14.23\%，经营利润率=5.33\%，预测 5 年），系统输出以下估值结果：

\begin{table}[H]
\centering
\caption{DCF 估值结果}
\begin{tabular}{@{}lr@{}}
\toprule
\textbf{指标} & \textbf{数值} \\
\midrule
WACC（加权平均资本成本）     & 10.00\% \\
企业价值（Enterprise Value） & 3,929,036,401.44 元 \\
权益价值（Equity Value）     & 3,939,589,827.81 元 \\
每股价值                      & 4.74 元 \\
预测年数                      & 5 年 \\
\bottomrule
\end{tabular}
\end{table}

\subsection{敏感性分析矩阵}

敏感性分析生成了 $9 \times 7$ 的每股价值矩阵，WACC 范围 8\%--12\%，永续增长率范围 1.0\%--4.0\%，帮助投资者理解估值的参数敏感度。

\subsection{AI 估值叙述}

系统成功调用 MiniMax LLM 生成约 4,700 字的专业估值分析报告，涵盖公司概况、关键指标、WACC 假设分析、现金流预测、终端价值方法论、估值结论及风险提示等内容。

\subsection{端到端性能}

\begin{table}[H]
\centering
\caption{系统性能指标}
\begin{tabular}{@{}lr@{}}
\toprule
\textbf{环节} & \textbf{耗时} \\
\midrule
PDF 上传与文本提取  & $\sim$5 秒 \\
LLM 数据提取       & $\sim$45 秒 \\
DCF 计算           & $<$ 1 秒 \\
敏感性分析         & $<$ 1 秒 \\
AI 叙述生成        & $\sim$50 秒 \\
\midrule
\textbf{全流程总计} & $\sim$100 秒 \\
\bottomrule
\end{tabular}
\end{table}

% ══════════════════════════════════════════
\section{技术挑战与解决方案}

\subsection{挑战一：openai/httpx 版本不兼容}

\textbf{问题}：\texttt{openai==1.51.0} 在创建 \texttt{httpx.AsyncClient} 时传递了 \texttt{proxies} 参数，而 \texttt{httpx>=0.28} 已移除该参数，导致 \texttt{AsyncClient.\_\_init\_\_() got an unexpected keyword argument 'proxies'} 错误。

\textbf{解决}：将 openai 升级至 \texttt{>=1.58.0,<2.0.0}（实际安装 1.109.1），该版本兼容 httpx 0.28+。

\subsection{挑战二：LLM 推理标签导致 JSON 解析失败}

\textbf{问题}：MiniMax-M2.7-highspeed 是推理型模型，在每次响应前自动输出 \texttt{<think>...</think>} 思考过程。当 \texttt{max\_tokens} 不足时，模型将全部 token 用于推理，最终 JSON 未能输出。

\textbf{解决}：
\begin{enumerate}[leftmargin=2em]
  \item 将 \texttt{max\_tokens} 从 2,000 提升至 16,000，为思考与输出留出充足空间；
  \item 编写正则表达式处理器，自动剥离 \texttt{<think>} 标签后提取 JSON 对象。
\end{enumerate}

\subsection{挑战三：PDF 关键数据分散}

\textbf{问题}：226 页年报的关键财务数据分散在不同章节（第 7--9 页的主要指标、第 84 页起的资产负债表、利润表等），直接提取前 12,000 字符无法覆盖核心数据。

\textbf{解决}：实现基于关键词权重评分的智能页面选择算法，优先提取包含``营业收入''``净利润''``资产负债表''等 17 个关键词的高价值页面，总提取量控制在 30,000 字以内。

\subsection{挑战四：环境变量加载路径}

\textbf{问题}：\texttt{config.py} 中 \texttt{load\_dotenv()} 默认在工作目录查找 \texttt{.env} 文件，但实际文件位于 \texttt{backend/.env} 子目录。

\textbf{解决}：显式指定路径 \texttt{load\_dotenv(\_backend\_dir / ".env")}，确保从项目根目录启动时正确加载配置。

% ══════════════════════════════════════════
\section{总结与展望}

\subsection{项目成果}

本项目成功构建了一套完整的 DCF 估值分析智能代理系统，实现了以下成果：

\begin{itemize}[leftmargin=2em]
  \item 实现了从 PDF 年报上传到 AI 自动提取 21 项财务指标的全自动化流程；
  \item 构建了完整的 FCFF-DCF 估值模型引擎，支持 5 年投射与 Gordon 永续增长终端价值；
  \item 生成 $9 \times 7$ 的 WACC-增长率敏感性矩阵，以热力图可视化呈现；
  \item 利用 MiniMax LLM 自动生成 4,000+ 字的专业估值叙述报告；
  \item 打造了美观的全栈 Web 应用，支持中英文双语与深色/浅色主题。
\end{itemize}

\subsection{未来改进方向}

\begin{enumerate}[leftmargin=2em]
  \item \textbf{多期数据分析}：支持多年财报对比，自动追踪关键指标趋势。
  \item \textbf{行业对标}：引入同行业可比公司数据，提供相对估值参考。
  \item \textbf{实时市场数据}：接入股票行情 API，自动获取当前股价与 Beta 值。
  \item \textbf{蒙特卡洛模拟}：基于参数分布进行随机模拟，输出估值概率分布。
  \item \textbf{报告导出}：支持将估值结果导出为 PDF/Excel 格式的专业报告。
  \item \textbf{多模型支持}：适配更多 LLM（如 GPT-4o、Claude 等），对比不同模型的提取准确度。
\end{enumerate}

% ══════════════════════════════════════════
\appendix
\section{核心数据模型定义}

以下为系统核心 Pydantic 数据模型，定义了前后端之间的数据契约：

\begin{lstlisting}[language=Python,caption=FinancialData 模型]
class FinancialData(BaseModel):
    company_name: str = ""
    ticker: Optional[str] = None
    currency: str = "CNY"
    fiscal_year: int = 2025
    revenue: float = 0
    revenue_growth: float = 0
    operating_income: float = 0
    operating_margin: float = 0.15
    net_income: float = 0
    depreciation_amortization: float = 0
    capital_expenditure: float = 0
    change_in_working_capital: float = 0
    total_debt: float = 0
    cash_and_equivalents: float = 0
    shares_outstanding: float = 0
    tax_rate: float = 0.25
    beta: float = 1.0
    risk_free_rate: float = 0.03
    market_return: float = 0.09
    cost_of_debt: float = 0.05
    current_stock_price: Optional[float] = None
\end{lstlisting}

\begin{lstlisting}[language=Python,caption=DCFParameters 模型]
class DCFParameters(BaseModel):
    projection_years: int = 5
    revenue_growth_rate: float = 0.08
    terminal_growth_rate: float = 0.025
    operating_margin: float = 0.15
    tax_rate: float = 0.25
    capex_ratio: float = 0.05
    da_ratio: float = 0.04
    nwc_ratio: float = 0.02
    wacc: float = 0.1
    discount_rate_override: Optional[float] = None
\end{lstlisting}

\begin{lstlisting}[language=Python,caption=DCFResult 模型]
class DCFResult(BaseModel):
    projections: list[FCFProjection]
    terminal_value: float
    pv_terminal_value: float
    pv_fcf_sum: float
    enterprise_value: float
    equity_value: float
    per_share_value: float
    current_price: Optional[float] = None
    upside_downside: Optional[float] = None
    wacc_used: float
    terminal_growth_used: float
\end{lstlisting}

\end{document}
